% !TEX program = xelatex
% Example: Complete thesis template with frontmatter, mainmatter, appendix, backmatter
% Demonstrates the standard structure of a university thesis using report class.
% Compile with: xelatex main.tex (twice for cross-references)

\documentclass[12pt, a4paper]{book}

% Essential packages
\usepackage{fontspec}
\usepackage{polyglossia}
\setmainlanguage[locale=mashriq]{arabic}
\setotherlanguage{english}
\newfontfamily\arabicfont[Script=Arabic]{Amiri}
\newfontfamily\englishfont{Latin Modern Roman}

\usepackage{amsmath, amssymb}
\usepackage{graphicx}
\usepackage{booktabs}
\usepackage{hyperref}
\usepackage[backend=biber, style=authoryear]{biblatex}
\usepackage{imakeidx}
\makeindex

% Dummy bibliography entries
\usepackage{filecontents}
\begin{filecontents*}{refs.bib}
@book{knuth1984,
  author = {Donald E. Knuth},
  title = {The TeXbook},
  publisher = {Addison-Wesley},
  year = {1984}
}
@book{lamport1994,
  author = {Leslie Lamport},
  title = {LaTeX: A Document Preparation System},
  publisher = {Addison-Wesley},
  year = {1994}
}
\end{filecontents*}
\addbibresource{refs.bib}

\begin{document}

%========== FRONTMATTER ==========
\frontmatter

\chapter*{صفحةُ العنوان}
\begin{center}
\vspace{2cm}
{\Large \textbf{عنوانُ الأطروحةِ هنا}}\\[1cm]
{\large اسمُ الباحث}\\[0.5cm]
{\large المُشرف: د. اسمُ المُشرف}\\[0.5cm]
{\large جامعةُ ... - قسمُ ...}\\[0.5cm]
{\large 2025}
\end{center}

\clearpage

\chapter*{الإهدار}
\begin{center}
إلى عائلتي وأصدقائي...
\end{center}

\clearpage

\chapter*{الشكرُ والتقدير}
أشكرُ مُشرفي على دعمهِ وتوجيهاتهِ القيّمة. كما أشكرُ زملاءَ البحثِ على التعاون.

\clearpage

\chapter*{الملخص}
هذا الملخصُ يَلخّصُ المشكلةَ والمنهجَ والنتائجَ في 150 إلى 350 كلمة. يُوضعُ هنا نصُّ الملخص.

\clearpage

\tableofcontents
\listoffigures
\listoftables

%========== MAINMATTER ==========
\mainmatter

\chapter{المُقدّمة}
\section{خلفيةُ البحث}
هذا الفصلُ يُقدّمُ خلفيةَ البحث. نُشيرُ إلى المرجعِ \textcite{knuth1984}.

\section{أهدافُ البحث}
الأهدافُ الرئيسةُ لهذا البحث...

\chapter{مراجعةُ الأدبيات}
\section{الدراساتُ السابقة}
مراجعةُ الأبحاثِ السابقةِ في المجال. نُشيرُ إلى \textcite{lamport1994}.

\chapter{المنهجية}
\section{تصميمُ البحث}
وصفُ المنهجِ المُتّبع.

\chapter{النتائج}
\section{النتائجُ الرئيسة}
عرضُ النتائجِ معَ الأشكالِ والجداول.

\begin{figure}[htbp]
\centering
\fbox{\parbox{8cm}{\centering [Figure placeholder]}}
\caption{نتيجةٌ تجريبية}
\label{fig:result}
\end{figure}

\begin{table}[htbp]
\centering
\begin{tabular}{lc}
\toprule
المُتغيّر & القيمة \\
\midrule
الدقة & 95.2\% \\
الزمن & 1.3 ثانية \\
\bottomrule
\end{tabular}
\caption{نتائجُ الأداء}
\label{tab:results}
\end{table}

\chapter{المناقشة}
\section{تفسيرُ النتائج}
كما نَرى في الشكلِ \ref{fig:result} والجدولِ \ref{tab:results}، النتائجُ تَدعمُ الفرضية.

\chapter{الخاتمة}
\section{الخلاصةُ والتوصيات}
خلاصةُ البحثِ والتوصياتُ المستقبلية.

%========== APPENDIX ==========
\appendix

\chapter{بياناتٌ إضافية}
هذا الملحقُ يَحتوي على بياناتٍ تكميلية.

\chapter{شفراتٌ برمجية}
هذا الملحقُ يَحتوي على الشفراتِ المُستخدمة.

%========== BACKMATTER ==========
\backmatter

\printbibliography
\printindex

\end{document}
